% ============ 26 Weil 猜想 (原书 pp.151-153) ============
\chapter{Weil 猜想}\label{sec:26}
\index{Weil 猜想}
\index{$\zeta$ 函数}
\index{Riemann 猜想}
\index{函数方程}

设 $X_0$ 是 $\mathbb{F}_q$ 上的非奇异射影簇.对每个 $m$，设 $N_m$ 是 $X_0$ 上坐标在 $\mathbb{F}_{q^m}$ 中的点的个数，并把 $X_0$ 的 $\zeta$ 函数定义为
\[
Z(X_0, t) = \exp\left( \sum_{m \geq 1} N_m \frac{t^{m}}{m} \right)
\]
\[
= 1 + \sum_{m \geq 1} N_m \frac{t^{m}}{m} + \frac{1}{2!}\left( \sum_{m \geq 1} N_m \frac{t^{m}}{m} \right)^{2} + \cdots
\]
它是系数在 $\mathbb{Q}$ 中的形式幂级数，即 $Z(X_0, t) \in \mathbb{Q}[[t]]$.注意
\[
\frac{d}{dt} \log Z(X_0, t) = \sum_{m \geq 1} N_m t^{m-1}.
\]
于是 $\frac{d}{dt} \log Z(X_0, t)$ 本质上是序列 $N_1, N_2, N_3, \ldots$ 的生成函数.

除记号的微小改动外，下面是 Weil 的\emph{Numbers of solutions of equations in finite fields}, Bull. AMS 55 (1949), 497--508 的逐字引用.

\begin{quote}
这以及其他我们这里不能讨论的例子，似乎对如下猜想性论断提供了一些支持.这些论断已知对曲线成立，但迄今为止我未能对高维簇证明.

设 $X_0$ 是定义在 $q$ 个元素的有限域 $\mathbb{F}_q$ 上的，维数 $d$ 的无奇异点簇\footnote{在这篇论文中，Weil 似乎用``簇''指``射影簇''.}.设 $N_m$ 是 $X_0$ 在次数为 $m$ 的 $\mathbb{F}_q$ 的扩张 $\mathbb{F}_{q^m}$ 上的有理点个数.则我们有
\[
\sum_{1}^{\infty} N_m t^{m-1} = \frac{d}{dt} \log Z(t),
\]
其中 $Z(t)$ 是 $t$ 的有理函数，满足函数方程
\[
Z\left( \frac{1}{q^{d}t} \right) = \pm q^{d\chi/2} t^{\chi} Z(t),
\]
其中 $\chi$ 等于 $X$ 的 Euler-Poincar\'e 示性数($X$ 的对角与其自身在乘积 $X \times X$ 中的交数).

此外，我们有：
\[
Z(t) = \frac{P_1(t)P_3(t) \cdots P_{2d-1}(t)}{P_0(t)P_2(t) \cdots P_{2d}(t)},
\]
其中 $P_0(t) = 1 - t$，$P_{2d}(t) = 1 - q^{d}t$，且对 $1 \leq r \leq 2d - 1$：
\[
P_r(t) = \prod_{i=1}^{\beta_r} (1 - \alpha_{r,i} t),
\]
其中 $\alpha_{r,i}$ 是绝对值为 $q^{r/2}$ 的代数整数.

最后，让我们把多项式 $P_r(t)$ 的次数 $\beta_r$ 称为簇 $X$ 的 Betti 数；于是 Euler-Poincar\'e 示性数 $\chi$ 由通常的公式 $\chi = \sum_r (-1)^{r}\beta_r$ 表示.手头的证据似乎表明，若 $Y$ 是定义在代数数域 $K$ 上的无奇异点簇，则由 $Y$ 模一个素理想 $\mathfrak{p}$ 约化得到的簇 $Y_{(\mathfrak{p})}$ 的 Betti 数，对所有至多有限个素理想 $\mathfrak{p}$，都等于 $Y$(看作复数域上的簇)在组合拓扑意义下的 Betti 数.
\end{quote}

课程的余下部分致力于证明这些引人注目的猜想.

条件 $|\alpha_{r,i}| = q^{r/2}$ 通常称为 Riemann 猜想.

尽管 Weil 在他 1949 年的论文中\footnote{或我所知道的任何其他地方，尽管文献中有许多作者反复作相反的陈述.可能 Weil 知道在特征 $p$ 中不存在带 $\mathbb{Q}$ 系数的好上同调理论(因为存在超奇异椭圆曲线)，因而不愿意猜想存在带特征零域系数的好上同调理论(当前术语中的``Weil 上同调'').回忆 Weil 小心地区分猜想与推测.}没有明确这样说，但显然他的猜想的形式部分地由代数拓扑的形式体系提示(见 Weil, \emph{Œuvres}, I, p.568).特别地，函数方程由 Poincar\'e 对偶提示，$\zeta$ 函数作为有理函数的形式由 Lefschetz 不动点公式提示.

\begin{remark}\label{rem:26.1}
设 $Y$ 是 $\operatorname{Spec} \mathbb{Z}$ 上有限型的概形.$Y$ 的闭点 $y$ 处的剩余域有限---设 $N(y)$ 是它的阶.$Y$ 的 $\zeta$ 函数定义为
\[
\zeta(Y, s) = \prod_{y} \frac{1}{1 - N(y)^{-s}},
\]
其中 $y$ 遍历 $Y$ 的闭点.该乘积收敛，并把 $\zeta(Y, s)$ 定义为 $\Re(s) > \dim Y$ 上的全纯函数.例如，若 $Y = \operatorname{Spec} \mathbb{Z}$，则 $\zeta(Y, s)$ 是 Riemann 的原始 $\zeta$ 函数.$\mathbb{F}_q$ 上的簇 $X_0$ 可以通过
\[
X_0 \to \operatorname{Spec} \mathbb{F}_q \to \operatorname{Spec} \mathbb{Z}
\]
(第二个映射由 $\mathbb{Z} \to \mathbb{Z}/p\mathbb{Z} \hookrightarrow \mathbb{F}_q$ 定义)看作 $\operatorname{Spec} \mathbb{Z}$ 上有限型的概形，我们将在下一节((4), p.158)看到
\[
Z(X_0, t) = \prod_{x} \frac{1}{1 - t^{\deg x}},
\]
从而
\[
\zeta(X_0, s) = Z(X_0, q^{-s}).
\]
因此，$X_0$ 的 Riemann 猜想说 $\zeta(X_0, s)$ 的极点都在直线 $\Re(s) = 0, 1, 2, \ldots, \dim X$ 上，零点都在直线 $\Re(s) = \frac{1}{2}, \frac{3}{2}, \ldots, \dim X - \frac{1}{2}$ 上.与原始 Riemann 猜想的类比是显然的.
\end{remark}
